\documentclass[book.tex]{subfiles}
\begin{document}

\chapter{Linear Systems and Matrix Algebra}
\label{chap:linear_systems}
\noindent\rule{11cm}{0.4pt} \\
\textbf{Prerequisites:} \autoref{chap:vectors}, \autoref{chap:optimization} \\
\textbf{Difficulty Level:} * \\
\noindent\rule{11cm}{0.4pt}
\vspace{5mm}

This chapter will introduce you to matrix notation and will introduce some common types of matrices including identity matrices, diagonal matrices, symmetric matrices, triangular matrices and tridiagonal matrices. We will then cover how to represent a system of linear algebraic equations in matrix form and how such systems can be solved through matrix inversion.

\section{Motivation}

Linear system of equations need to be solved in numerous applications including parabolic interpolation where the minima at every iteration is determined by fitting a parabola through three points and finding the minima. The general form of a parabola is  $f(x) = ax^2+bx+c$, and the minima is obtained by equating the derivative $f'(x)$ to zero, which implies that the extrema occurs at $\dfrac{-b}{2a}$.  Using the constraints that $x_1$, $x_2$ and $x_3$ lie on a parabola, we can write the following system of equations to solve for $a$, $b$ and $c$ to obtain the parabola and thereby the minima.

\begin{align}
ax_1^2 + bx_1 + c &= f(x_1)\\
ax_2^2 + bx_2 + c &= f(x_2)\\
ax_3^2 + bx_3 + c &= f(x_3)
\end{align}
\[	
\underbrace{
	\begin{bmatrix}
	x_1^2 & x_1 & 1 \\
	x_2^2 & x_2 & 1 \\
	x_3^2 & x_3 & 1 \\
	\end{bmatrix}}_{\text{Known}}
\underbrace{
	\begin{bmatrix}
	a\\
	b\\
	c
	\end{bmatrix}}_{\text{Unknown}} = 
\underbrace{
	\begin{bmatrix}
	f(x_1)\\
	f(x_2)\\
	f(x_3)
	\end{bmatrix}}_{\text{Known}}
\]

\section{Matrix Algebra Overview}

A \textit{matrix} consists of a rectangular array of elements represented by a single symbol. As depicted in Figure \ref{fig:1}, $A$ is the shorthand notation for the matrix and $a_{ij}$ designates an individual \textit{element} of the matrix.

\begin{marginfigure}
	\centering
	% \includegraphics[width = 2.7in]{figures/part1b/linear_systems/matrixindices.png}
 \includegraphics[width = 2.7in]{figures/part1b/linear_systems/matrixindices_0.png}
	\centering
	\caption{A matrix indicating the indices of each element}
	\label{fig:1}
\end{marginfigure}

A horizontal set of elements is called a \textit{row} and a vertical set is called a \textit{column}. The first subscript i always designates the number of the row in which the element lies. The second subscript j designates the column. For example, element $a_{23}$ is in row 2 and column 3. 

The matrix in the Figure \ref{fig:1} has \textit{m} rows and \textit{n} columns and is said to have a dimension of \textit{m} by \textit{n} (or \textit{m} $\times$ \textit{n}). It is referred to as an \textit{m} by \textit{n} matrix.

\section{Common Types of Matrices}
\subsection{Row and Column Matrices}
Matrices with row dimension $m$ = 1 are called \textit{row vectors}. Similarly, matrices with column dimension $n$ = 1 are called \textit{column vectors}. A representation of a row vector is given below.
\begin{align}
	b = [b_1 \hspace{10pt }b_2  \hspace{10pt} \cdots  \hspace{10pt} b_n]
\end{align}

\subsection{Square Matrices}
Matrices where $m$ = $n$ are called \textit{square matrices}. The diagonal consisting of the elements $a_{11}$, $a_{22}$, and $a_{33}$ is termed the \textit{principal} or \textit{main diagonal} of the matrix. 

\begin{margintable}
	\[
	A=
	\begin{bmatrix}
	a_{11} & a_{12} & a_{13} \\
	a_{21} & a_{22} & a_{23} \\
	a_{31} & a_{32} & a_{33}
	\end{bmatrix}
	\]
\end{margintable}

Square matrices are particularly important when solving sets of simultaneous linear equations. For such systems, the number of equations (corresponding to rows) and the number of unknowns (corresponding to columns) must be equal for a unique solution to be possible.


There are a number of special forms of square matrices that are important and should
be noted:

\subsection{Symmetric Matrix}
\begin{margintable}
	\[
	S=
	\begin{bmatrix}
	a_{11} & a_{12} & a_{13} \\
	a_{12} & a_{22} & a_{23} \\
	a_{13} & a_{23} & a_{33}
	\end{bmatrix}
	\]
\end{margintable}
A \textit{symmetric matrix} is one where the rows equal the columns - that is, $a_{ij}$= $a_{ji}$ for all $i$'s and $j$'s.
\vspace{-13pt}

\subsection{Diagonal Matrix}
\begin{margintable}
	\[
	D=
	\begin{bmatrix}
	a_{11} &  &  \\
	& a_{22} &  \\
	&  & a_{33}
	\end{bmatrix}
	\]
\end{margintable}
A \textit{diagonal matrix} is a square matrix where all elements off the main diagonal are equal to zero. Note that where large blocks of elements are zero, they are left blank.

\subsection{Identity Matrix}
An \textit{identity matrix} is a diagonal matrix where all elements on the main diagonal are equal to 1. The identity matrix has properties similar to unity. That is,
\begin{margintable}
	\[
	I=
	\begin{bmatrix}
	1 &  &  \\
	& 1 &  \\
	&  & 1
	\end{bmatrix}
	\]
\end{margintable}
\vspace{-20pt}
\begin{align}
AI = IA = A
\end{align}

\subsection{Upper Triangular Matrix}
An \textit{upper triangular matrix} is one where all the elements below the main diagonal are zero.
\begin{margintable}
	\[
	U=
	\begin{bmatrix}
	a_{11} & a_{12} & a_{13} \\
	& a_{22} & a_{23} \\
	&  & a_{33}
	\end{bmatrix}
	\]
\end{margintable}
\vspace{-8pt}

\subsection{Lower Triangular Matrix}
\vspace{-6pt}
\begin{margintable}
	\[
	L=
	\begin{bmatrix}
	a_{11} &  &  \\
	a_{21} & a_{22} &  \\
	a_{31} & a_{32} & a_{33}
	\end{bmatrix}
	\]
\end{margintable}
A \textit{lower triangular matrix} is one where all elements above the main diagonal are zero.

\subsection{Banded Matrix}

\begin{margintable}
	\[
	B=
	\begin{bmatrix}
	a_{11} & a_{12} & &  \\
	a_{21} & a_{22} & a_{23} & \\
	& a_{32} & a_{33} & a_{34} \\
	&  & a_{43} & a_{44} \\
	\end{bmatrix}
	\]
\end{margintable}

A \textit{banded matrix} has all elements equal to zero, with the exception of a band centered on the main diagonal. The matrix given to the right has a bandwidth of 3 and is given a special name - the \textit{tridiagonal matrix}.

\subsection{Transpose Matrix}
\begin{margintable}
	\[
	A'=
	\begin{bmatrix}
	a_{11} & a_{21} & a_{31} \\
	a_{12} & a_{22} & a_{32} \\
	a_{13} & a_{23} & a_{33}
	\end{bmatrix}
	\]
\end{margintable}

The \textit{transpose} of a matrix involves transforming its rows into columns and its columns into rows. In other words, the element $a_{ij}$ of the transpose is equal to the $a_{ji}$ element of the original matrix. The transpose of [A] is given to the right.

\subsection{Permutation Matrix}
\begin{margintable}
	\[
	P=
	\begin{bmatrix}
	0 & 0  & 1 \\
	0 & 1 & 0 \\
	1 & 0 & 0
	\end{bmatrix}
	\]
\end{margintable}

A \textit{permutation} matrix (also called a \textit{transposition} matrix) is an identity matrix with rows and columns interchanged. For example, a permutation matrix that is constructed by switching the first and third rows and columns of a 3 $\times$ 3 identity matrix is displayed to the right. Left multiplying a matrix $A$ by this matrix, as in $PA$, will switch the corresponding rows of $A$. Right multiplying, as in $AP$, will switch the corresponding columns.

\section{Matrix Operating Rules}
\subsection{Matrix Addition}
Addition of two matrices, say, $A$ and $B$, is accomplished by adding corresponding
terms in each matrix. The elements of the resulting matrix $C$ are computed as\\
\begin{align}
c_{ij} = a_{ij} + b_{ij}
\end{align}
for $i$ = 1, 2, . . . ,$m$ and $j$ = 1, 2, . . . ,$ n$. 
The addition of matrices satisfies the \textit{commutative} and \textit{associative} properties:

\begin{margintable}
\textbf{Matrix Addition Identities:}
\begin{Verbatim}
 [A] + [B] = [B] + [A] 
([A] + [B]) + [C] = [A] + ([B] + [C])
\end{Verbatim}
\end{margintable}
\subsection{Matrix Multiplication}
The multiplication of a matrix $A$ by a scalar $g$ is obtained by multiplying every element of $A$ by $g$.

\begin{marginfigure}
	\centering
	% \includegraphics[width = 2.7in]{figures/part1b/linear_systems/matrixmultiplication.PNG}	\centering
 \includegraphics[width = 2.7in]{figures/part1b/linear_systems/matrixmultiplication_0.png}	\centering
	\caption{Visual Depiction of Matrix Multiplication}
	\label{fig:2}
\end{marginfigure}
The product of two matrices is represented as $C = AB$, where the elements of $C$ are defined as
\begin{align}
c_{ij} = \sum_{k=1}^{n}a_{ik}b_{kj}
\end{align}

where \textit{n} = the column dimension of $A$ and the row dimension of $B$. Figure \ref{fig:2} depicts how the rows and columns line up in matrix multiplication. According to this definition, matrix multiplication can be performed only if the first matrix has as many columns as the number of rows in the second matrix. 

\begin{margintable}
	\textbf{Matrix Multiplication Identities:}
\begin{Verbatim}
([A][B])[C] = [A]([B][C])
[A]([B] + [C]) = [A][B] + [A][C]
\end{Verbatim}
\end{margintable}

If the dimensions of the matrices are suitable, matrix multiplication is \textit{associative} and \textit{distributive}. However, multiplication is not generally \textit{commutative}. That is, the order of matrix multiplication is important.

\begin{margintable}
\begin{lstlisting}
[A][B] $\neq$ [B][A]
\end{lstlisting}
\end{margintable}

Although multiplication is possible, matrix division is not a defined operation. However, if a matrix $A$ is square and nonsingular, there is another matrix $A^{-1}$, called the inverse of $A$, for which

\begin{align}
AA^{-1} = A^{-1}A = I
\end{align}


The inverse of a 2 $\times$ 2 matrix can be represented by
\begin{align}
A^{-1} = \dfrac{1}{a_{11}a_{22}-a_{12}a_{21}}	\begin{bmatrix}
a_{22} & -a_{12} \\
-a_{21} & a_{11} 
\end{bmatrix}
\end{align}

\begin{margintable}
	\begin{verbatim}
	Example:
	>> A = [3 5 1;5 8 1;-7 -2 4]
	A =
	3     5     1
	5     8     1
	-7    -2     4
	\end{verbatim}
\end{margintable}

\begin{margintable}
	\begin{verbatim}	
	>> det(A)
	ans =	
	
	\end{verbatim}
\end{margintable}

\begin{margintable}
	\begin{verbatim}
	>> transpose(A) 
	ans =
	3     5    -7
	5     8    -2
	1     1     4
	\end{verbatim}
\end{margintable}

\begin{margintable}
	\begin{verbatim}	
	>> inv(A)	
	ans =	
	2.6154   -1.6923   -0.2308
	-2.0769    1.4615    0.1538
	3.5385   -2.2308   -0.0769
	\end{verbatim}
\end{margintable}

Let's consider a matrix $A$ as given to the right and calculate its \textit{Determinant}, \textit{Transpose} and \textit{Inverse}. %The determinant of a matrix in Matlab is evaluated by the function \texttt{det}. The transpose of of a matrix in Matlab is given by the function \texttt{transpose}. The inverse of of a matrix in Matlab is calculated by the function \texttt{inv}. The command window session of these function calls is given to the right.

\section{Representing Linear Equations as Matrices}
It can be observed that matrices provide a concise notation for representing simultaneous linear equations. For example, a 3 $\times$ 3 set of linear equations,
\begin{align}
a_{11}x_1 + a_{12}x_2 + a_{13}x_3 = b_1\\
a_{21}x_1 + a_{22}x_2 + a_{23}x_3 = b_2\\
a_{31}x_1 + a_{32}x_2 + a_{33}x_3 = b_3
\end{align}

can be expressed as 
\begin{align}
Ax = b
\end{align}

where $A$ is the matrix of coefficients, $b$ is the column vector of constants and $x$ is the column vector of unknowns.

A formal way to obtain a solution using matrix algebra is to multiply each side of the equation by the inverse of $A$ to yield
\begin{align}
x = A^{-1}b
\end{align}

\begin{margintable}
	\textbf{Matrix Representation of Linear Equations:}
	\\
	\vspace{20pt}

\[	
	\underbrace{
	\begin{bmatrix}
	a_{11} & a_{12} & a_{13} \\
	a_{12} & a_{22} & a_{23} \\
	a_{13} & a_{23} & a_{33}
	\end{bmatrix}}_{\text{A}}
	\underbrace{
	\begin{bmatrix}
			x_{1}\\
			x_{2}\\
			x_{3}
	\end{bmatrix}}_{\text{x}} = 
	\underbrace{
	\begin{bmatrix}
	b_{1}\\
	b_{2}\\
	b_{3}
	\end{bmatrix}}_{\text{b}}s
\]
		
	
\end{margintable}

It should be noted that the above equation is true only for systems with equal number of equations and unknowns, that is \textit{m} = \textit{n}. Systems with more equations (rows) than unknowns (columns), \textit{m} > \textit{n}, are said to be \textit{over-determined}. A typical example is least-squares regression where an equation with \textit{n} coefficients is fit to \textit{m} data points. Conversely, systems with less equations than unknowns, \textit{m} < \textit{n}, are said to be \textit{under-determined}. A typical example of under-determined systems is numerical optimization.
%\section{Solving Linear Equations with Matlab}
%MATLAB provides two direct ways to solve systems of linear algebraic equations. 

%\begin{enumerate}
%\item The most efficient way is to employ the backslash, or %\textit{left-division operator} as \texttt{A$\textbackslash$b }
%\item
%The second way is to use matrix inversion as \texttt{x = inv(A) $\times$ b}
%\end{enumerate}

%\textbf{Example:} Solve the following example using the two methods discussed above.
%\begin{align}
%	3x_1 + 2x_2 = 123\\
%-4x_1 +  3x_3 = 652\\
%	 5x_2 +7x_3 = 234
%\end{align}
%\end{document}
%\section{Exercises}
%\begin{enumerate}
%    \item \textcolor{red}{TODO}
%\end{enumerate}
\end{document}